% --- Archon physics formalization source begin ---
% archon:physics
% archon:covers PhyXMiniProblems/problem_phyx_mini_0850.lean
% archon:source-report reports/phyx_mini/problem_phyx_mini_0850.source.json

\chapter{Physics problem phyx\_mini\_0850}
\label{ch:PhyXMiniProblems_problem_phyx_mini_0850}

\paragraph{Problem source.}
\#\# Physical scenario

The figure shows four sides of a \textbackslash{}( 3.0\textbackslash{},\textbackslash{}mathrm\{cm\} 	imes 3.0\textbackslash{},\textbackslash{}mathrm\{cm\} 	imes 3.0\textbackslash{},\textbackslash{}mathrm\{cm\} \textbackslash{}) cube.

\#\# Question

What is the net flux through these four sides?

\#\# Answer choices

- A: A: 80N m\textasciicircum{}2/C

- B: B: 10N m\textasciicircum{}2/C

- C: C: 20N m\textasciicircum{}2/C

- D: D: 0N m\textasciicircum{}2/C

\#\# Figure caption (auxiliary; use the image as primary evidence)

The image shows the top view of a 3.0 cm × 3.0 cm × 3.0 cm cube with several elements:

1. **Electric Field Lines**: There are multiple magenta arrows representing electric field lines entering and exiting the square plane of the cube. These lines are parallel and uniformly spaced.

2. **Labeling**: The lines are labeled with circled numbers (1 through 4), indicating specific positions of the field lines.

3. **Text**: 
   - Top left: "Top view of a 3.0 cm × 3.0 cm × 3.0 cm cube"
   - Center of the cube: "E" with an arrow over it, indicating it is a vector.
   - Near the lower right field line: "500 N/C," indicating the magnitude of the electric field.

4. **Angle**: An angle of 30° is noted between one of the field lines and a horizontal line, indicating the inclination of these lines.

This visual suggests a representation of the electric field through and around the cube.

\#\# Dataset metadata

Electrostatics; Physical Model Grounding Reasoning, Spatial Relation Reasoning

\paragraph{Recorded answer/context.}
D: D: 0N m\textasciicircum{}2/C

\paragraph{Figure/image path.}
/root/proposal\_for\_physic/science-mango/phyx\_mini\_run/phyx\_data/test\_image/850.png

\paragraph{Formalization target.}
create a compiling Lean file with sorry bodies at `PhyXMiniProblems/problem\_phyx\_mini\_0850.lean`. The Lean declarations must preserve the physical quantities, dimensions or dimensional roles, figure labels, governing-law hypotheses, and final relation expressed by this problem.
Use Mathlib/Physlib names found through LeanExplore where available. If a physics API is missing, introduce faithful local abstractions rather than scalar placeholder aliases.

\begin{theorem}[Physics formalization target]
\label{thm:physics:phyx_mini_0850:target}
\lean{PhyXMiniProblems.ProblemPhyXMini0850.problem_phyx_mini_0850}
\uses{def:physics:phyx-mini-0850:phyxminiproblems-problemphyxmini0850-electricfluxinnewtonsquaremeterspercoulomb, def:physics:phyx-mini-0850:phyxminiproblems-problemphyxmini0850-uniformfieldcubesetup, def:physics:phyx-mini-0850:phyxminiproblems-problemphyxmini0850-matchessuppliedtopviewfigure, def:physics:phyx-mini-0850:phyxminiproblems-problemphyxmini0850-isphysicalfoursidecubegeometry, def:physics:phyx-mini-0850:phyxminiproblems-problemphyxmini0850-isuniformtopviewelectricfield, def:physics:phyx-mini-0850:phyxminiproblems-problemphyxmini0850-satisfiesfoursideplanarfluxlaws, def:physics:phyx-mini-0850:phyxminiproblems-problemphyxmini0850-recordeddatasetanswer, def:physics:phyx-mini-0850:phyxminiproblems-problemphyxmini0850-answermatchesnetfoursideflux}
The assigned autoformalize agent should translate this physics problem into Lean declarations in the covered file, with theorem and lemma proof bodies written as `by sorry`.
\end{theorem}
\begin{proof}
This is an autoformalization task, not a proof task. Produce statements that can later be proved by the physics prover without weakening the physical model.
\end{proof}

% --- Archon declaration topology begin ---
\section{Declaration topology}
\begin{definition}[electric Field Strength Dimension]
  \label{def:physics:phyx-mini-0850:phyxminiproblems-problemphyxmini0850-electricfieldstrengthdimension}
  \lean{PhyXMiniProblems.ProblemPhyXMini0850.electricFieldStrengthDimension}
  The dimension M L T⁻² C⁻¹ of electric-field strength (N/C)
\end{definition}
\begin{proof}
  This is the declared model definition of electric Field Strength Dimension.
\end{proof}
\begin{definition}[area Dimension]
  \label{def:physics:phyx-mini-0850:phyxminiproblems-problemphyxmini0850-areadimension}
  \lean{PhyXMiniProblems.ProblemPhyXMini0850.areaDimension}
  The physical dimension L² of one square side of the cube
\end{definition}
\begin{proof}
  This is the declared model definition of area Dimension.
\end{proof}
\begin{definition}[electric Flux Dimension]
  \label{def:physics:phyx-mini-0850:phyxminiproblems-problemphyxmini0850-electricfluxdimension}
  \lean{PhyXMiniProblems.ProblemPhyXMini0850.electricFluxDimension}
  \uses{def:physics:phyx-mini-0850:phyxminiproblems-problemphyxmini0850-electricfieldstrengthdimension, def:physics:phyx-mini-0850:phyxminiproblems-problemphyxmini0850-areadimension}
  Electric flux has dimension M L³ T⁻² C⁻¹, equivalently N m²/C in coherent SI units
\end{definition}
\begin{proof}
  This is the declared model definition of electric Flux Dimension.
\end{proof}
\begin{definition}[Length Quantity]
  \label{def:physics:phyx-mini-0850:phyxminiproblems-problemphyxmini0850-lengthquantity}
  \lean{PhyXMiniProblems.ProblemPhyXMini0850.LengthQuantity}
  A nonnegative, unit-independent physical length
\end{definition}
\begin{proof}
  This is the declared model definition of Length Quantity.
\end{proof}
\begin{definition}[Area Quantity]
  \label{def:physics:phyx-mini-0850:phyxminiproblems-problemphyxmini0850-areaquantity}
  \lean{PhyXMiniProblems.ProblemPhyXMini0850.AreaQuantity}
  \uses{def:physics:phyx-mini-0850:phyxminiproblems-problemphyxmini0850-areadimension}
  A nonnegative, unit-independent physical area
\end{definition}
\begin{proof}
  This is the declared model definition of Area Quantity.
\end{proof}
\begin{definition}[Electric Field Strength Quantity]
  \label{def:physics:phyx-mini-0850:phyxminiproblems-problemphyxmini0850-electricfieldstrengthquantity}
  \lean{PhyXMiniProblems.ProblemPhyXMini0850.ElectricFieldStrengthQuantity}
  \uses{def:physics:phyx-mini-0850:phyxminiproblems-problemphyxmini0850-electricfieldstrengthdimension}
  A nonnegative, unit-independent electric-field magnitude
\end{definition}
\begin{proof}
  This is the declared model definition of Electric Field Strength Quantity.
\end{proof}
\begin{definition}[Electric Flux Quantity]
  \label{def:physics:phyx-mini-0850:phyxminiproblems-problemphyxmini0850-electricfluxquantity}
  \lean{PhyXMiniProblems.ProblemPhyXMini0850.ElectricFluxQuantity}
  \uses{def:physics:phyx-mini-0850:phyxminiproblems-problemphyxmini0850-electricfluxdimension}
  A signed, unit-independent electric flux. Its sign is determined by the outward orientation of the relevant cube side
\end{definition}
\begin{proof}
  This is the declared model definition of Electric Flux Quantity.
\end{proof}
\begin{definition}[length In Meters]
  \label{def:physics:phyx-mini-0850:phyxminiproblems-problemphyxmini0850-lengthinmeters}
  \lean{PhyXMiniProblems.ProblemPhyXMini0850.lengthInMeters}
  \uses{def:physics:phyx-mini-0850:phyxminiproblems-problemphyxmini0850-lengthquantity}
  Read a physical length in metres
\end{definition}
\begin{proof}
  This is the declared model definition of length In Meters.
\end{proof}
\begin{definition}[length In Centimeters]
  \label{def:physics:phyx-mini-0850:phyxminiproblems-problemphyxmini0850-lengthincentimeters}
  \lean{PhyXMiniProblems.ProblemPhyXMini0850.lengthInCentimeters}
  \uses{def:physics:phyx-mini-0850:phyxminiproblems-problemphyxmini0850-lengthquantity}
  Read a physical length in centimetres
\end{definition}
\begin{proof}
  This is the declared model definition of length In Centimeters.
\end{proof}
\begin{definition}[area In Square Meters]
  \label{def:physics:phyx-mini-0850:phyxminiproblems-problemphyxmini0850-areainsquaremeters}
  \lean{PhyXMiniProblems.ProblemPhyXMini0850.areaInSquareMeters}
  \uses{def:physics:phyx-mini-0850:phyxminiproblems-problemphyxmini0850-areaquantity}
  Read a physical area in square metres
\end{definition}
\begin{proof}
  This is the declared model definition of area In Square Meters.
\end{proof}
\begin{definition}[electric Field Strength In Newtons Per Coulomb]
  \label{def:physics:phyx-mini-0850:phyxminiproblems-problemphyxmini0850-electricfieldstrengthinnewtonspercoulomb}
  \lean{PhyXMiniProblems.ProblemPhyXMini0850.electricFieldStrengthInNewtonsPerCoulomb}
  \uses{def:physics:phyx-mini-0850:phyxminiproblems-problemphyxmini0850-electricfieldstrengthquantity}
  Read electric-field strength in newtons per coulomb
\end{definition}
\begin{proof}
  This is the declared model definition of electric Field Strength In Newtons Per Coulomb.
\end{proof}
\begin{definition}[electric Flux In Newton Square Meters Per Coulomb]
  \label{def:physics:phyx-mini-0850:phyxminiproblems-problemphyxmini0850-electricfluxinnewtonsquaremeterspercoulomb}
  \lean{PhyXMiniProblems.ProblemPhyXMini0850.electricFluxInNewtonSquareMetersPerCoulomb}
  \uses{def:physics:phyx-mini-0850:phyxminiproblems-problemphyxmini0850-electricfluxquantity}
  Read signed electric flux in newton-square-metres per coulomb
\end{definition}
\begin{proof}
  This is the declared model definition of electric Flux In Newton Square Meters Per Coulomb.
\end{proof}
\begin{definition}[Lateral Face]
  \label{def:physics:phyx-mini-0850:phyxminiproblems-problemphyxmini0850-lateralface}
  \lean{PhyXMiniProblems.ProblemPhyXMini0850.LateralFace}
  The four vertical cube faces visible in the top view. The constructors retain the circled labels 1 through 4 printed around the square in the raster
\end{definition}
\begin{proof}
  This is the declared model definition of Lateral Face.
\end{proof}
\begin{definition}[Top View Side]
  \label{def:physics:phyx-mini-0850:phyxminiproblems-problemphyxmini0850-topviewside}
  \lean{PhyXMiniProblems.ProblemPhyXMini0850.TopViewSide}
  The callouts identify side 1 as the left face, side 2 as the upper face, side 3 as the right face, and side 4 as the lower face in the top view
\end{definition}
\begin{proof}
  This is the declared model definition of Top View Side.
\end{proof}
\begin{definition}[figure Side Placement]
  \label{def:physics:phyx-mini-0850:phyxminiproblems-problemphyxmini0850-figuresideplacement}
  \lean{PhyXMiniProblems.ProblemPhyXMini0850.figureSidePlacement}
  \uses{def:physics:phyx-mini-0850:phyxminiproblems-problemphyxmini0850-lateralface, def:physics:phyx-mini-0850:phyxminiproblems-problemphyxmini0850-topviewside}
  Placement of each numbered callout in the supplied top view
\end{definition}
\begin{proof}
  This is the declared model definition of figure Side Placement.
\end{proof}
\begin{definition}[figure Outward Unit Normal]
  \label{def:physics:phyx-mini-0850:phyxminiproblems-problemphyxmini0850-figureoutwardunitnormal}
  \lean{PhyXMiniProblems.ProblemPhyXMini0850.figureOutwardUnitNormal}
  \uses{def:physics:phyx-mini-0850:phyxminiproblems-problemphyxmini0850-lateralface}
  Dimensionless outward unit normals in top-view coordinates (x,y,z). The third coordinate is perpendicular to the page
\end{definition}
\begin{proof}
  This is the declared model definition of figure Outward Unit Normal.
\end{proof}
\begin{definition}[degrees]
  \label{def:physics:phyx-mini-0850:phyxminiproblems-problemphyxmini0850-degrees}
  \lean{PhyXMiniProblems.ProblemPhyXMini0850.degrees}
  Convert a numerical degree readout into a physical angle
\end{definition}
\begin{proof}
  This is the declared model definition of degrees.
\end{proof}
\begin{definition}[top View Direction]
  \label{def:physics:phyx-mini-0850:phyxminiproblems-problemphyxmini0850-topviewdirection}
  \lean{PhyXMiniProblems.ProblemPhyXMini0850.topViewDirection}
  The dimensionless direction in the top-view plane making angle with the positive horizontal axis
\end{definition}
\begin{proof}
  This is the declared model definition of top View Direction.
\end{proof}
\begin{definition}[Cube Geometry]
  \label{def:physics:phyx-mini-0850:phyxminiproblems-problemphyxmini0850-cubegeometry}
  \lean{PhyXMiniProblems.ProblemPhyXMini0850.CubeGeometry}
  \uses{def:physics:phyx-mini-0850:phyxminiproblems-problemphyxmini0850-lengthquantity, def:physics:phyx-mini-0850:phyxminiproblems-problemphyxmini0850-areaquantity, def:physics:phyx-mini-0850:phyxminiproblems-problemphyxmini0850-lateralface}
  The cubical geometry relevant to the requested four lateral faces. The side regions and sample points keep the flux law tied to actual subsets and points of Physlib's three-dimensional space
\end{definition}
\begin{proof}
  This is the declared model definition of Cube Geometry.
\end{proof}
\begin{definition}[Supplied Top View Figure]
  \label{def:physics:phyx-mini-0850:phyxminiproblems-problemphyxmini0850-suppliedtopviewfigure}
  \lean{PhyXMiniProblems.ProblemPhyXMini0850.SuppliedTopViewFigure}
  \uses{def:physics:phyx-mini-0850:phyxminiproblems-problemphyxmini0850-lengthquantity, def:physics:phyx-mini-0850:phyxminiproblems-problemphyxmini0850-electricfieldstrengthquantity, def:physics:phyx-mini-0850:phyxminiproblems-problemphyxmini0850-lateralface, def:physics:phyx-mini-0850:phyxminiproblems-problemphyxmini0850-topviewside}
  Typed evidence transcribed from the primary raster. The length and field labels are physical quantities rather than bare scalar aliases
\end{definition}
\begin{proof}
  This is the declared model definition of Supplied Top View Figure.
\end{proof}
\begin{definition}[Uniform Field Cube Setup]
  \label{def:physics:phyx-mini-0850:phyxminiproblems-problemphyxmini0850-uniformfieldcubesetup}
  \lean{PhyXMiniProblems.ProblemPhyXMini0850.UniformFieldCubeSetup}
  \uses{def:physics:phyx-mini-0850:phyxminiproblems-problemphyxmini0850-electricfieldstrengthquantity, def:physics:phyx-mini-0850:phyxminiproblems-problemphyxmini0850-electricfluxquantity, def:physics:phyx-mini-0850:phyxminiproblems-problemphyxmini0850-lateralface, def:physics:phyx-mini-0850:phyxminiproblems-problemphyxmini0850-cubegeometry, def:physics:phyx-mini-0850:phyxminiproblems-problemphyxmini0850-suppliedtopviewfigure}
  The physical setup and its independent observables. electricField is the full Physlib field. Its vector components below are interpreted as coherent SI N/C readouts, while electricFieldStrength supplies the corresponding dimensionful magnitude. Neither fluxThrough nor netFourSideFlux is defined from an answer choice
\end{definition}
\begin{proof}
  This is the declared model definition of Uniform Field Cube Setup.
\end{proof}
\begin{definition}[Matches Supplied Top View Figure]
  \label{def:physics:phyx-mini-0850:phyxminiproblems-problemphyxmini0850-matchessuppliedtopviewfigure}
  \lean{PhyXMiniProblems.ProblemPhyXMini0850.MatchesSuppliedTopViewFigure}
  \uses{def:physics:phyx-mini-0850:phyxminiproblems-problemphyxmini0850-lengthincentimeters, def:physics:phyx-mini-0850:phyxminiproblems-problemphyxmini0850-electricfieldstrengthinnewtonspercoulomb, def:physics:phyx-mini-0850:phyxminiproblems-problemphyxmini0850-lateralface, def:physics:phyx-mini-0850:phyxminiproblems-problemphyxmini0850-figuresideplacement, def:physics:phyx-mini-0850:phyxminiproblems-problemphyxmini0850-degrees, def:physics:phyx-mini-0850:phyxminiproblems-problemphyxmini0850-uniformfieldcubesetup}
  Facts read from 850.png: a 3.0 cm cube, four numbered side callouts, a uniform-looking family of parallel arrows labeled 500 N/C, and a 30° angle. The image contains no numerical flux readout, so no requested answer is introduced here
\end{definition}
\begin{proof}
  This is the declared model definition of Matches Supplied Top View Figure.
\end{proof}
\begin{definition}[Is Physical Four Side Cube Geometry]
  \label{def:physics:phyx-mini-0850:phyxminiproblems-problemphyxmini0850-isphysicalfoursidecubegeometry}
  \lean{PhyXMiniProblems.ProblemPhyXMini0850.IsPhysicalFourSideCubeGeometry}
  \uses{def:physics:phyx-mini-0850:phyxminiproblems-problemphyxmini0850-lengthinmeters, def:physics:phyx-mini-0850:phyxminiproblems-problemphyxmini0850-areainsquaremeters, def:physics:phyx-mini-0850:phyxminiproblems-problemphyxmini0850-lateralface, def:physics:phyx-mini-0850:phyxminiproblems-problemphyxmini0850-figureoutwardunitnormal, def:physics:phyx-mini-0850:phyxminiproblems-problemphyxmini0850-uniformfieldcubesetup}
  Geometric facts about the four congruent square faces. Their common area is the square of the positive cube edge; each representative point lies on its face; and the outward normals agree with the numbered top-view callouts
\end{definition}
\begin{proof}
  This is the declared model definition of Is Physical Four Side Cube Geometry.
\end{proof}
\begin{definition}[Is Uniform Top View Electric Field]
  \label{def:physics:phyx-mini-0850:phyxminiproblems-problemphyxmini0850-isuniformtopviewelectricfield}
  \lean{PhyXMiniProblems.ProblemPhyXMini0850.IsUniformTopViewElectricField}
  \uses{def:physics:phyx-mini-0850:phyxminiproblems-problemphyxmini0850-electricfieldstrengthinnewtonspercoulomb, def:physics:phyx-mini-0850:phyxminiproblems-problemphyxmini0850-topviewdirection, def:physics:phyx-mini-0850:phyxminiproblems-problemphyxmini0850-uniformfieldcubesetup}
  The field is spatially uniform at the observation time and has the magnitude and in-plane direction recorded by the setup. This is a governing property of the depicted field, not a statement about electric flux
\end{definition}
\begin{proof}
  This is the declared model definition of Is Uniform Top View Electric Field.
\end{proof}
\begin{definition}[Satisfies Four Side Planar Flux Laws]
  \label{def:physics:phyx-mini-0850:phyxminiproblems-problemphyxmini0850-satisfiesfoursideplanarfluxlaws}
  \lean{PhyXMiniProblems.ProblemPhyXMini0850.SatisfiesFourSidePlanarFluxLaws}
  \uses{def:physics:phyx-mini-0850:phyxminiproblems-problemphyxmini0850-areainsquaremeters, def:physics:phyx-mini-0850:phyxminiproblems-problemphyxmini0850-electricfluxinnewtonsquaremeterspercoulomb, def:physics:phyx-mini-0850:phyxminiproblems-problemphyxmini0850-lateralface, def:physics:phyx-mini-0850:phyxminiproblems-problemphyxmini0850-uniformfieldcubesetup}
  For a constant field on a planar side, oriented flux is A * inner(E,n\_out). The second field is additivity over the four-side boundary partition. Both are general physical laws and neither asserts the requested zero value
\end{definition}
\begin{proof}
  This is the declared model definition of Satisfies Four Side Planar Flux Laws.
\end{proof}
\begin{lemma}[opposite Side Flux Contributions cancel]
  \label{lem:physics:phyx-mini-0850:phyxminiproblems-problemphyxmini0850-oppositesidefluxcontributions-cancel}
  \lean{PhyXMiniProblems.ProblemPhyXMini0850.oppositeSideFluxContributions_cancel}
  \uses{def:physics:phyx-mini-0850:phyxminiproblems-problemphyxmini0850-electricfluxinnewtonsquaremeterspercoulomb, def:physics:phyx-mini-0850:phyxminiproblems-problemphyxmini0850-uniformfieldcubesetup, def:physics:phyx-mini-0850:phyxminiproblems-problemphyxmini0850-isphysicalfoursidecubegeometry, def:physics:phyx-mini-0850:phyxminiproblems-problemphyxmini0850-isuniformtopviewelectricfield, def:physics:phyx-mini-0850:phyxminiproblems-problemphyxmini0850-satisfiesfoursideplanarfluxlaws}
  Opposite side normals are negatives of one another. Uniformity and the planar-face flux law therefore make the side 1/3 and side 2/4 contributions cancel pairwise
\end{lemma}
\begin{proof}
  The conclusion follows from the governing relations and figure readouts named in the statement of opposite Side Flux Contributions cancel.
\end{proof}
\begin{definition}[Answer Choice]
  \label{def:physics:phyx-mini-0850:phyxminiproblems-problemphyxmini0850-answerchoice}
  \lean{PhyXMiniProblems.ProblemPhyXMini0850.AnswerChoice}
  Labels of the four electric-flux choices printed in the source
\end{definition}
\begin{proof}
  This is the declared model definition of Answer Choice.
\end{proof}
\begin{definition}[answer Flux In Newton Square Meters Per Coulomb]
  \label{def:physics:phyx-mini-0850:phyxminiproblems-problemphyxmini0850-answerfluxinnewtonsquaremeterspercoulomb}
  \lean{PhyXMiniProblems.ProblemPhyXMini0850.answerFluxInNewtonSquareMetersPerCoulomb}
  \uses{def:physics:phyx-mini-0850:phyxminiproblems-problemphyxmini0850-answerchoice}
  Flux readout in N m²/C printed beside each answer label
\end{definition}
\begin{proof}
  This is the declared model definition of answer Flux In Newton Square Meters Per Coulomb.
\end{proof}
\begin{definition}[recorded Dataset Answer]
  \label{def:physics:phyx-mini-0850:phyxminiproblems-problemphyxmini0850-recordeddatasetanswer}
  \lean{PhyXMiniProblems.ProblemPhyXMini0850.recordedDatasetAnswer}
  \uses{def:physics:phyx-mini-0850:phyxminiproblems-problemphyxmini0850-answerchoice}
  The answer label recorded by the source dataset
\end{definition}
\begin{proof}
  This is the declared model definition of recorded Dataset Answer.
\end{proof}
\begin{definition}[Answer Matches Net Four Side Flux]
  \label{def:physics:phyx-mini-0850:phyxminiproblems-problemphyxmini0850-answermatchesnetfoursideflux}
  \lean{PhyXMiniProblems.ProblemPhyXMini0850.AnswerMatchesNetFourSideFlux}
  \uses{def:physics:phyx-mini-0850:phyxminiproblems-problemphyxmini0850-electricfluxinnewtonsquaremeterspercoulomb, def:physics:phyx-mini-0850:phyxminiproblems-problemphyxmini0850-uniformfieldcubesetup, def:physics:phyx-mini-0850:phyxminiproblems-problemphyxmini0850-answerchoice, def:physics:phyx-mini-0850:phyxminiproblems-problemphyxmini0850-answerfluxinnewtonsquaremeterspercoulomb}
  A displayed choice agrees with the independently modeled net four-side flux. This relation does not define the net flux from the choice
\end{definition}
\begin{proof}
  This is the declared model definition of Answer Matches Net Four Side Flux.
\end{proof}
% --- Archon declaration topology end ---
% --- Archon physics formalization source end ---
